\chapter{Apr.~9 --- Classification of Curves, Part 2}

\section{Ramification}

\begin{remark}
  As before, let $f : X \to Y$
  be a surjective morphism of
  (smooth projective) curves.
  We will try to relate $\Omega_{X / Y}$
  to ``ramification.''
\end{remark}

\begin{definition}
  Fix $P \in X$ and set $Q = f(P)$.
  Choose uniformizers $u \in \OO_{X, P}$ and
  $v \in \OO_{Y, Q}$. Now
  $f^* v = c u^{e_P}$ for some
  $c \in \OO_{X, P}^\times$ and
  $e_P \in \Z_{\ge 1}$ (note that
  $v$ vanishes at $Q$, so $f^* v$ must vanish
  at $P$, hence $e_P \ne 0$). Equivalently,
  $e_P = \mathrm{coeff}_P(f^* Q)$. We say
  that $f$ is \emph{ramified} at $P$
  if $e_P > 1$.
\end{definition}

\begin{example}
  Let $\Affine^1 \to \Affine^1$ be given
  by $x \mapsto x^d$. Then
  \[
    e_P =
    \begin{cases}
      d & \text{if } P = 0, \\
      1 & \text{if } P \ne 0.
    \end{cases}
  \]
\end{example}

\begin{definition}
  The \emph{ramification divisor} of $X$ is
  \[
    \Ram_f
    = \sum_{P \in X} (e_P - 1) P.
  \]
  We call $\Supp(\Ram_f)$ the
  \emph{ramification locus}.
\end{definition}

\begin{remark}
  Assume going forward that $f$ is separable.
  Note that $(\Omega_X)_P = \OO_{X, P} du$
  since
  \[
    (\Omega_X)_{(P)}
    \cong \m_P / \m_P^2.
  \]
  Similarly, we have
  $(f^* \Omega_Y)_P \cong (\Omega_Y)_Q \otimes \OO_{X, P}$,
  generated by $dv \otimes 1$, and
  $(\OO_{\Ram_f})_P \cong \OO_{X, P} / u^{e_P - 1} \OO_{X, P}$.
  Now we have a short exact sequence
  \[
    \begin{tikzcd}[row sep=tiny]
      0 \ar[r] & (f^* \Omega_Y)_P
      \ar[r] & (\Omega_X)_P
      \ar[r] & (\Omega_{X / Y})_P \ar[r] & 0 \\
             & dv \otimes 1\ar[r, mapsto]
             & d(cu^{e_P})
    \end{tikzcd}
  \]
  where $d(cu^{e_P}) = u^{e_P} dc + cu^{e_P - 1} du = cu^{e_P - 1} du$ since $c$ is a unit.
  Thus $\Omega_{X / Y} = \OO_{\Ram_f}$, so we have
  \[
    \begin{tikzcd}
      0 \ar[r]
      & f^* \Omega_Y \ar[r]
      & \Omega_X \ar[r] & \OO_{\Ram_f} \ar[r] & 0
    \end{tikzcd}
  \]
  Applying $- \otimes \Omega_X^\vee$, we
  get a short exact sequence
  \[
    \begin{tikzcd}
      0 \ar[r] & \Omega_X^\vee \otimes f^* \Omega_Y
      \ar[r] & \OO_X \ar[r] & \OO_{\Ram_f} \ar[r] & 0
    \end{tikzcd}
  \]
  Note that $\OO_{\Ram_f}$ is supported
  at points, so
  $\Omega_X^\vee \otimes \OO_{\Ram_f} \cong \OO_{\Ram_f}$.
  Thus
  $\Omega_X^\vee
    \otimes f^* \Omega_Y
    \cong \OO_X(-{\Ram_f})$,
  so
  \[
    \Omega_X
    \cong f^* \Omega_Y \otimes \OO_X(\Ram_f).
  \]
  This implies the following result:
\end{remark}

\begin{theorem}[Riemann-Hurwitz formula]
  We have the following:
  \begin{enumerate}
    \item $\Ram_f + f^* K_Y \sim K_X$,
    \item $2 g(X) - 2 = (\deg f)(2g(Y) - 2) + \deg \Ram_f$.
  \end{enumerate}
\end{theorem}

\begin{remark}
  We have the following:
  \begin{enumerate}
    \item $g(X) \ge g(Y)$,
    \item $g(X) = g(Y)$ occurs only when
      $d = 1$ and there is no ramification, or
      $X, Y$ are either both genus $0$ or
      both genus $1$.
  \end{enumerate}
\end{remark}

\section{Classification of Higher Genus Curves}

\begin{remark}
  Morally, hyperelliptic curves are
  ``curves with 2-to-1 maps to $\PP^1$.''
\end{remark}

\begin{definition}
  The \emph{gonality} of a curve $X$, denoted
  $\gon(X)$, is the
  smallest degree of an invertible sheaf
  $\mathcal{L}$ on $X$ with
  $h^0(X, \mathcal{L}) \ge 2$.
\end{definition}

\begin{example}
  We have the following:
  \begin{enumerate}
    \item $\gon(\PP^1) = 1$ (computed
      by $\OO_{\PP^1}(1)$, or equivalently,
      $\OO_{\PP^1}(P)$ for any $P \in \PP^1$).
    \item $\gon(E) = 2$ for an
      elliptic curve $E$, since
        $\deg \mathcal{L}
        = h^0(E, \mathcal{L})
        - h^0(E, \mathcal{L}^\vee)$,
        thus $\deg \mathcal{L} = 1$
        implies $h^0(E, \mathcal{L}) = 1$
        and $\deg \mathcal{L} = 2$ implies $h^0(E, \mathcal{L}) = 2$.
  \end{enumerate}
\end{example}

\begin{lemma}
  The gonality $\gon(X)$ is the smallest
  degree of a finite map $X \to \PP^1$.
\end{lemma}

\begin{proof}
  If $f : X \to \PP^1$ is a degree $d$ map,
  then $\mathcal{L} = f^* \OO_{\PP^1}(1)$
  has degree $d$ and
  $h^0(X, \mathcal{L}) \ge 2$, so
  \[\deg f \ge \gon(X).\]
  Conversely, fix a minimal degree invertible
  sheaf $\mathcal{L}$ on $X$ with
  $h^0(X, \mathcal{L}) \ge 2$. Then for $P \in X$,
  \[
    h^0(X, \mathcal{L}(-P))
    = h^0(X, \mathcal{L}) - 1,
  \]
  by minimality, and in fact
  $h^0(X, \mathcal{L}) = 2$ by a similar
  argument. Thus
  $\mathcal{L}$ is globally generated and
  induces a map $f : X \to \PP^1$
  with $f^* \OO_{\PP^1}(1) \cong \mathcal{L}$.
  Thus $\deg \mathcal{L} = (\deg f)(\deg \OO_{\PP^1}(1)) = \deg f$.
\end{proof}

\begin{remark}
  An elliptic curve has a degree $2$
  map $E \to \PP^1$, so by the
  Riemann-Hurwitz formula,
  \[
    \deg \Ram_f = 4.
  \]
  So there are 4 distinct points of
  ramification. In fact, the
  4 ramification points (up to isomorphisms
  of $\PP^1$) classifies the
  elliptic curve up to isomorphism.
\end{remark}

\begin{definition}
  A curve $X$ is \emph{hyperelliptic} if
  $\gon(X) = 2$ (i.e. the minimal degree
  map $X \to \PP^1$ has degree $2$).
\end{definition}

\begin{remark}
  A hyperelliptic curve $X$ with a
  degree $2$ morphism $X \to \PP^1$ can be
  constructed from the data of $\Ram_f$.
\end{remark}

\begin{example}
  Every smooth genus $2$ curve is
  hyperelliptic and has a 2-to-1 map
  $X \to \PP^1$ ramified at 6 points.

  To see this, note that if
  $g = 2$, then $\deg \omega_X = 2 \cdot 2 - 2 = 2$.
  Now
  \[
    h^0(X, \omega_X)
    = h^1(X, \omega_X)
    + \deg \omega_X + 1 - 2
    = h^0(X, \OO_X) + 1 = 2,
  \]
  so $\gon(X) \le 2$. But
  $\gon(X) = 1$ would imply $X \cong \PP^1$,
  so $\gon(X) = 2$.
\end{example}

\begin{remark}
  For $g \ge 2$, we have
  that $\omega_X$ is globally
  generated.

  To see this, it suffices to show that
  $h^0(X, \omega_X) - h^0(X, \omega_X(-P)) = 1$
  for all $P \in X$. We have
  \[
    h^0(X, \omega_X(-P))
    = h^1(X, \OO_X(P))
    = h^0(X, \OO_X(P))
    - \deg P + g - 1,
  \]
  where the first equality is
  by Serre duality and the second equality
  is by Riemann-Roch. Now we must have
  $h^0(X, \OO_X(P)) = 1$ as otherwise
  $h^0(X, \OO_X(P)) \ge 2$ would imply
  $\gon(X) = 1$. Thus
  \[
    h^0(X, \omega_X(-P))
    = g - 1
    = h^1(X, \OO_X) - 1
    = h^0(X, \omega_X) - 1,
  \]
  where the last equality is by
  Serre duality. Thus $\omega_X$ is
  globally generated.
\end{remark}

\begin{remark}
  For $g \ge 2$, when is $\omega_X$ very ample?
  Note that $\omega_X$ is not very ample
  if and only if there exist
  $P, Q \in X$ such that
  \[
    h^0(X, \omega_X)
    - h^0(X, \omega_X(-P - Q))
    \ne 2.
  \]
  Since we already know $\omega_X$ is
  globally generated, this is equivalent
  to saying that
  \[
    h^0(X, \omega_X)
    - h^0(X, \omega_X(-P - Q))
    = 1.
  \]
  By Serre duality, this equivalent to there
  being $P, Q \in X$ such that
  \[
    g - h^1(X, \OO_X(P + Q)) = 1,
  \]
  which by Riemann-Roch is equivalent to
  there being $P, Q \in X$ such that
  \[
    g - (h^0(X, \OO_X(P + Q)) - \deg(P + Q) + g - 1) = 1,
  \]
  which is equivalent to
  $h^0(X, \OO_X(P + Q)) = 2$. This is
  then equivalent to $\gon(X) = 2$.
  This proves:
\end{remark}

\begin{theorem}
  For a curve $X$ with
  $g \ge 2$, either $X$ is
  hyperelliptic or $\omega_X$ is very ample.
\end{theorem}

\begin{definition}
  A \emph{canonical curve} $X$ is a curve
  of genus $\ge 2$ with $\omega_X$ very ample.
\end{definition}

\begin{remark}
  If $X$ is a canonical curve, then there
  exists an embedding
  $X \hookrightarrow \PP^{g - 1}$ induced by
  $\omega_X$.
\end{remark}

\begin{example}
  The following are some
  classifications of canonical
  curves:
  \begin{enumerate}
    \item $g = 3$: quartic curve in $\PP^2$,
    \item $g = 4$: intersection of
      a quadric with a cubic in $\PP^3$,
    \item $g = 5$: intersection of
      $3$ quadrics in $\PP^4$,
    \item $g = 6$: intersection of
      $G(2, 5)$ in $\PP^9$ with a linear
      subspace $\PP^5 \subseteq \PP^9$ and
      a quadric.
  \end{enumerate}
\end{example}
